\documentclass[11pt]{article}
\usepackage{icmlko}
\begin{document}

\title{한 번 잰 수는 부호도 못 맞춘다 --- 그리고 내가 미리 적은 규칙이 규약과 어긋났다}
\author{월드모델 실험실}
\date{노트 358 · 2026년 8월 1일}
\maketitle

\begin{abstract}
노트 357이 팔 넷을 \emph{한 번씩} 적합해 표로 적었고, 그중 하나(등급 필터를
A$\cdot$B $\to$ A\textasciitilde{}E 로 풀기)를 ``짝으로 재야 배포 규칙에
걸 수 있다''로 남겼다. 쟀다. \textbf{한 번 적합의 도메인별 차는 열 중
\emph{넷}에서 부호가 틀렸다.} 제일 중요한 자리에서 틀렸다 --- 대표 수치인
팝업이 한 번 적합에서 $-0.0095$ 였는데 짝 열둘에서
\textbf{$+0.0140$(9/12)} 다. 아이돌은 $-0.1000$ 이 $+0.0026$ 이 된다.
\textbf{노트 352가 ``수준은 못 잰다''고 한 것의 실무 형태다} --- 한 번 적합의
도메인 칸은 근거가 아니라 숫자다.

그리고 \textbf{내가 이 실험을 위해 미리 적은 결정 규칙이 규약 275와
어긋나 있었다.} 나는 ``판이 9/12 미만이면 안 넣는다''로 적었는데, 275는
판이 미결정 폭 안이면 \emph{배포가 정한다}고 말한다. 두 조항이 정반대
답을 낸다. 275를 따랐다 --- 그것이 먼저 있고 더 넓고, 노트 351이 똑같은
자리에서 그렇게 했다. \textbf{판 $+0.0010$(6/12) · 팝업 $+0.0140$(9/12) ·
제일 나쁜 도메인 $|t|=0.07$ 이라 넣는다.} 새 기준값 판 \textbf{0.4862}
(분모 2,723).
\end{abstract}

\section{짝으로 재니 부호가 뒤집힌다}

\begin{center}\includegraphics[width=\linewidth]{figs/flip.pdf}\end{center}

\begin{center}
\begin{tabular}{lrrl}
\hline
도메인 & 한 번 적합 & 짝 열둘 & \\
\hline
\textbf{팝업}(대표) & $-0.0095$ & $\mathbf{+0.0140}$ (9/12) & \textbf{뒤집힘} \\
아이돌   & $-0.1000$ & $+0.0026$ (6/12) & \textbf{뒤집힘} \\
도서     & $+0.0286$ & $-0.0019$ (6/12) & \textbf{뒤집힘} \\
게임     & $-0.0062$ & $+0.0064$ (8/12) & \textbf{뒤집힘} \\
시장팝업 & $+0.0350$ & $+0.0015$ (7/12) & \\
펀딩     & $+0.0151$ & $+0.0028$ (5/12) & \\
웹툰     & $+0.0126$ & $+0.0007$ (6/12) & \\
모바일   & $+0.0024$ & $+0.0014$ (7/12) & \\
애니     & $-0.0005$ & $-0.0009$ (7/12) & \\
세계애니 & $-0.0006$ & $-0.0001$ (8/12) & \\
\hline
\end{tabular}
\end{center}

\textbf{열 중 넷이 부호가 반대이고, 크기는 더 나쁘다} --- 아이돌은 한 번
적합에서 $-0.1000$ 인데 짝으로는 $+0.0026$ 이다. 40배다. 남은 여섯도
짝 값이 한 번 값의 3분의 1 이하로 줄어든다(시장팝업 $0.0350 \to 0.0015$).

까닭은 노트 352가 이미 적었다 --- 유보 표집 오차가 한 번 적합에는 통째로
실리고 짝에는 두 팔에 똑같이 실려 상쇄된다. 얇은 도메인일수록 심하다
(아이돌 유보 51 · 팝업 65). \textbf{노트 357의 표에서 도메인 칸은 읽으면
안 되는 것이었다.} 그 글이 결론으로 쓴 것은 판과 기준별 $\rho$(둘 다
구간을 붙여 적었다)라 무사하지만, 표의 팝업 칸 $-0.0095$ 는 부호부터
틀렸다.

\textbf{규약: 도메인 칸은 짝으로만 적는다.} 한 번 적합의 도메인별 차는
표에 넣지 않는다 --- 넣으면 읽힌다.

\section{내 사전 등록이 규약과 어긋났다}

이 실험을 걸면서 결정 규칙을 미리 적었다.

\begin{quote}
판 $>0$ 이고 9/12 이상 · 팝업이 5/12 이상 $\to$ 넣는다 \\
팝업이 4/12 이하 $\to$ 안 넣는다 \\
\textbf{판이 9/12 미만 $\to$ 안 넣는다(바꿀 근거가 없다)}
\end{quote}

결과는 판 $+0.0010$(6/12) · 팝업 $+0.0140$(9/12)이다. 셋째 조항이 켜져서
``안 넣는다''가 된다.

\textbf{그런데 그 조항은 규약 275와 정반대다.} 275는 ``① 판이 미결정
폭($0.010$) 안이고 ② 어느 도메인도 $|t|>2.77$ 로 안 나빠지면 \emph{배포
쪽 이득이 정한다}''이다. 판이 미결정이라는 것은 275에서 \emph{유지하라}가
아니라 \emph{판 밖에서 정하라}는 뜻이라고 그 글이 명시해 뒀다. 내
셋째 조항은 같은 상황을 ``유지''로 읽었다.

여기서 조건을 확인하면 ① 판 $+0.0010 \ll 0.010$ --- 미결정 ② 제일 나쁜
도서가 $-0.0019/0.0273$ 로 $|t| = 0.07$ --- 문턱 밖이다. 그러면 배포가
정하고, 배포 수치가 $+0.0140$(9/12)이다.

\textbf{275를 따른다.} 셋이 근거다.
\begin{itemize}
\item 275가 먼저 있고 더 넓다 --- 노트별 사전 등록이 상시 규약을 조용히
덮으면 안 된다.
\item \textbf{노트 351이 똑같은 자리에서 그렇게 했다} --- 판
$+0.0002$(7/12) 미결정 · 팝업 $+0.0168$(11/12) $\to$ 넣었다. 여기서 안
넣으면 같은 상황을 두 번 다르게 판단하는 것이다.
\item \textbf{자료 결정의 근거가 점수 밖에 있다.} 노트 352가 시장팝업
유보 104행에서 \texttt{label\_trust} 등급이 정확도를 전혀 예측하지 않음을
쟀다 --- 깨끗한 B+C 가 $0.2147$ 로 전체 $0.2880$ 보다 낮고 같은 크기
무작위와 견줘 $p=0.907$ 이다. \textbf{다른 도메인에서, 팝업 점수를 보기
전에 나온 근거다.}
\end{itemize}

\textbf{노트 353의 규약대로 기제까지 적는다} --- 미리 적은 규칙을 뒤집을
때는 왜 그 규칙이 틀렸는지 보여야 한다. 여기서 틀린 것은 자료가 아니라
\emph{내가 쓴 셋째 조항}이고, 틀린 방식은 상시 규약과의 충돌이다.

\section{넣었다}

계수 방법은 안 건드렸다 --- 노트 357이 \texttt{organizer\_claim} 은 수준이
곱셈으로 부풀어 섞으면 순위가 크기가 아니라 기준으로 선다는 것을 보였다.
여기서 푼 것은 등급뿐이라 그 위험이 없다.

\begin{center}
\begin{tabular}{lrr}
\hline
 & 전(노트 357) & 지금 \\
\hline
팝업 행       & 75 (학습 16 · 유보 59) & \textbf{89} (학습 24 · 유보 65) \\
판           & 0.4814 & \textbf{0.4862} \\
판 분모      & 2,717 & 2,723 \\
판 95\%      & $[0.4485,\,0.5122]$ & $[0.4534,\,0.5162]$ \\
팝업         & 0.4905 & 0.4810 \\
팝업 95\%    & $[0.251,\,0.661]$ & $[0.223,\,0.661]$ \\
못 가리는 쌍 & 29 & 30 \\
\hline
\end{tabular}
\end{center}

\textbf{못 가리는 쌍이 하나 늘었다.} 아이돌이 $0.2137 \to 0.1137$ 로
내려가 시장팝업($0.1715$)과 겹쳤기 때문이다. 재는 힘이 는 처치가 아니다 ---
이건 배포 규칙이 넣으라고 한 처치다. 둘은 다른 목표이고 노트 352~356 이
넷 다 앞쪽이었던 것과 대조된다.

\textbf{그리고 학습이 16 $\to$ 24 다.} 노트 281이 확정한 잎 문턱은 22 이고
팝업은 그 아래에 있었다 --- 팝업 전용 축이 \emph{후보에도 못 오르는}
상태였다. 지금 처음으로 넘었다. 지금 챔피언에는 팝업 전용 축이 없어서
당장 쓸 데는 없지만, 노트 276이 만들고 노트 281이 ``문턱 아래라 죽는다''로
접은 넷(\texttt{pop\_comp} · \texttt{pop\_wkend} · \texttt{pop\_holid} ·
\texttt{pop\_stores})이 서랍에 있다.

\section{남은 것}

\begin{center}
\begin{tabular}{ll}
\hline
갈래 & 지금 아는 것 \\
\hline
\textbf{팝업 전용 축 넷} & 노트 276이 만들고 281이 접었다(학습 16 $<$ 문턱 22). \\
                        & 이제 24 다. \textbf{다음 순번} \\
기준별 라벨 정규화 & 노트 357. 부풀림이 사라지지만 과제가 바뀐다 \\
만화 2026 스무 건 & 노트 356. 열한 번째 채점 도메인 \\
만화가 왜 누르나 & 노트 356. 라벨 종류로는 안 갈린다 \\
\hline
\end{tabular}
\end{center}

첫 줄이 이 노트가 연 문이다. 노트 281이 ``팝업 전용 축은 순위로는 결재가
안 난다''고 접은 것이 학습 16행 때문이었고, 그 16이 24가 되었다.

\end{document}
